% !PRINTED 606
\weberpage{606}

Es handelt sich nun um den Grenzwert, dem diese Ausdrücke sich nähern, wenn $s=1$ wird.

Was die linke Seite betrifft, so folgt aus (17), daß
\weq{21}{\left\{\begin{aligned}
&\mathrm{Lim}\,(s-1)\,Q_1 = hg \text{ endlich und von Null verschieden,}\\
&\mathrm{Lim}\,Q_2\,Q_3\ldots Q_h \text{ endlich (vielleicht Null)}
\end{aligned}\right.}
ist. Zur Beurteilung der rechten Seite bemerken wir, daß die Primideale $\mathfrak{p}$ vom ersten oder vom zweiten Grade sind, und daß demnach $N(\mathfrak{p})$ gleich einer natürlichen Primzahl oder gleich dem Quadrat einer natürlichen Primzahl ist. Ferner stützen wir uns auf den bekannten Satz der Analysis, nach dem ein unendliches Produkt
\[
(1-q_1)\,(1-q_2)\,(1-q_3)\ldots
\]
unbedingt konvergiert und einen von Null verschiedenen Grenzwert hat, wenn
\[
q_1+q_2+q_3+\cdots
\]
eine unbedingt konvergente Reihe ist, und keines der $q_i=1$ ist.

Wir teilen danach die Primideale $\mathfrak{p}$ in zwei Arten, $\mathfrak{p}_1, \mathfrak{p}_2$, und zwar sollen die $\mathfrak{p}_1$ vom ersten Grade sein und zum Exponenten $\varphi=1$ gehören (also $\varepsilon=h$). Die $\mathfrak{p}_2$ sollen entweder vom zweiten Grade sein, oder zu einem höheren Exponenten $\varphi$ gehören. \gesperrt{Es sind also $\mathfrak{p}_1$ die zur Hauptklasse $\overline{A}_1$ gehörigen Primideale ersten Grades.}

Dann ist
\[
N(\mathfrak{p}_1)^{\varphi} = N(\mathfrak{p}_1) = p_1
\]
eine natürliche Primzahl und
\[
N(\mathfrak{p}_2)^{\varphi} = p_2^{k}
\]
mindestens die 2te (höchstens die $2h$te) Potenz einer Primzahl. Nach dem erwähnten analytischen Satze ist dann das Produkt
\[
\Pi\bigl[1-N(\mathfrak{p}_2)^{-s\varphi}\bigr]^{\varepsilon}
\]
für $s=1$ endlich und von Null verschieden, und es folgt aus (20), (21) der Satz:

6. \gesperrt{Der Grenzwert}
\[
\mathop{\mathrm{Lim}}\limits_{s=1} \; \Pi\,\dfrac{s-1}{\Pi\bigl[1-N(\mathfrak{p}_1)^{-s}\bigr]^{h}}
\]
\gesperrt{ist endlich, wenn $\mathfrak{p}_1$ die Primideale ersten Grades der Hauptklasse $\overline{A}_1$ durchläuft.}

Ein besonderes Interesse nimmt die Frage in Anspruch, ob dieser Grenzwert verschwinden kann.
